\documentclass{article}
\usepackage{hyperref}
\begin{document}
\section*{Stable destination and a changing actor after transmission}
\paragraph{Admitted causal cut.}
Both request policies send four bytes to the same native radio
receiver and the same duty. A prebound request names its chosen
actor at dispatch. The late form leaves only the choice of actor
to the authorized local dock; it does not grant another physical,
radio, or legal capability. The independent report is preserved
and the native ledgers are identical up to request reception.
The dock's simulated callback runs after the native sender and
receiver have been charged. Before actuation, the controller
injects the specified independent RWARE obstruction and applies
an explicit local rule. It retains the current actor if that cell
is open, otherwise it may regrant to the other open actor, or
withholds the action if regrant is not authorized. Except for
binding time, the two modes share the authority and occupancy
rules. This callback cut is a modeled cross-plant interleaving,
not a claim that native SubT or RWARE itself schedules a
distributed authority transfer.

\paragraph{Witness and comparator.}
In two runs the same old actor A is visible, both cells are open
at dispatch, and the same four bytes reach the dock. With no later
obstruction A still acts; with later interference on A and an
authorized local regrant, B performs the one actual movement.
An irrevocable sender choice based on the shared view cannot
serve both, even when its receiver rightly refuses a stale
physical binding. The late form rebinds only from the dock's
own trusted, current authority and physical access. If both cells
are blocked, or regrant is forbidden, it does not act.
A direct late check has the same results. This control shows
a cost of prebinding, not an advance over the strongest local
indirection under the same interfaces. The dock's continuing
address, permissions, visibility and actor actuation are explicit
premises, not inferred from connectivity.

\paragraph{Measured outcomes.}
Each arm has 192 matched cases, prebound versus late-bound versus
same-info direct late check. Actual settlements are 88, 110, 110;
exact result arrivals are 72, 90, 90. One outcome is not the
other: the dock may have performed the named action but the
native return is blocked by a competing receive debit. All arms
preserve the same 128 independent reports. Where prior side
traffic occupies the dock or caller ledger, the test keeps its order;
it also evaluates a separate declared case in which it can run
later. No universal early/late scheduling law or field delivery
bound is inferred. Lean checks the common-view obstruction,
validity of the later scoped choice and refusal when no current
actor is supplied. The fixed native adapter dependencies and
authorizations are not an independently chosen I1/I2 mission.
\end{document}
